% Options for packages loaded elsewhere
% Options for packages loaded elsewhere
\PassOptionsToPackage{unicode}{hyperref}
\PassOptionsToPackage{hyphens}{url}
\PassOptionsToPackage{dvipsnames,svgnames,x11names}{xcolor}
%
\documentclass[
  11pt,
]{article}
\usepackage{xcolor}
\usepackage[margin=1in]{geometry}
\usepackage{amsmath,amssymb}
\setcounter{secnumdepth}{-\maxdimen} % remove section numbering
\usepackage{iftex}
\ifPDFTeX
  \usepackage[T1]{fontenc}
  \usepackage[utf8]{inputenc}
  \usepackage{textcomp} % provide euro and other symbols
\else % if luatex or xetex
  \usepackage{unicode-math} % this also loads fontspec
  \defaultfontfeatures{Scale=MatchLowercase}
  \defaultfontfeatures[\rmfamily]{Ligatures=TeX,Scale=1}
\fi
\usepackage{lmodern}
\ifPDFTeX\else
  % xetex/luatex font selection
\fi
% Use upquote if available, for straight quotes in verbatim environments
\IfFileExists{upquote.sty}{\usepackage{upquote}}{}
\IfFileExists{microtype.sty}{% use microtype if available
  \usepackage[]{microtype}
  \UseMicrotypeSet[protrusion]{basicmath} % disable protrusion for tt fonts
}{}
\makeatletter
\@ifundefined{KOMAClassName}{% if non-KOMA class
  \IfFileExists{parskip.sty}{%
    \usepackage{parskip}
  }{% else
    \setlength{\parindent}{0pt}
    \setlength{\parskip}{6pt plus 2pt minus 1pt}}
}{% if KOMA class
  \KOMAoptions{parskip=half}}
\makeatother
% Make \paragraph and \subparagraph free-standing
\makeatletter
\ifx\paragraph\undefined\else
  \let\oldparagraph\paragraph
  \renewcommand{\paragraph}{
    \@ifstar
      \xxxParagraphStar
      \xxxParagraphNoStar
  }
  \newcommand{\xxxParagraphStar}[1]{\oldparagraph*{#1}\mbox{}}
  \newcommand{\xxxParagraphNoStar}[1]{\oldparagraph{#1}\mbox{}}
\fi
\ifx\subparagraph\undefined\else
  \let\oldsubparagraph\subparagraph
  \renewcommand{\subparagraph}{
    \@ifstar
      \xxxSubParagraphStar
      \xxxSubParagraphNoStar
  }
  \newcommand{\xxxSubParagraphStar}[1]{\oldsubparagraph*{#1}\mbox{}}
  \newcommand{\xxxSubParagraphNoStar}[1]{\oldsubparagraph{#1}\mbox{}}
\fi
\makeatother

\usepackage{color}
\usepackage{fancyvrb}
\newcommand{\VerbBar}{|}
\newcommand{\VERB}{\Verb[commandchars=\\\{\}]}
\DefineVerbatimEnvironment{Highlighting}{Verbatim}{commandchars=\\\{\}}
% Add ',fontsize=\small' for more characters per line
\usepackage{framed}
\definecolor{shadecolor}{RGB}{241,243,245}
\newenvironment{Shaded}{\begin{snugshade}}{\end{snugshade}}
\newcommand{\AlertTok}[1]{\textcolor[rgb]{0.68,0.00,0.00}{#1}}
\newcommand{\AnnotationTok}[1]{\textcolor[rgb]{0.37,0.37,0.37}{#1}}
\newcommand{\AttributeTok}[1]{\textcolor[rgb]{0.40,0.45,0.13}{#1}}
\newcommand{\BaseNTok}[1]{\textcolor[rgb]{0.68,0.00,0.00}{#1}}
\newcommand{\BuiltInTok}[1]{\textcolor[rgb]{0.00,0.23,0.31}{#1}}
\newcommand{\CharTok}[1]{\textcolor[rgb]{0.13,0.47,0.30}{#1}}
\newcommand{\CommentTok}[1]{\textcolor[rgb]{0.37,0.37,0.37}{#1}}
\newcommand{\CommentVarTok}[1]{\textcolor[rgb]{0.37,0.37,0.37}{\textit{#1}}}
\newcommand{\ConstantTok}[1]{\textcolor[rgb]{0.56,0.35,0.01}{#1}}
\newcommand{\ControlFlowTok}[1]{\textcolor[rgb]{0.00,0.23,0.31}{\textbf{#1}}}
\newcommand{\DataTypeTok}[1]{\textcolor[rgb]{0.68,0.00,0.00}{#1}}
\newcommand{\DecValTok}[1]{\textcolor[rgb]{0.68,0.00,0.00}{#1}}
\newcommand{\DocumentationTok}[1]{\textcolor[rgb]{0.37,0.37,0.37}{\textit{#1}}}
\newcommand{\ErrorTok}[1]{\textcolor[rgb]{0.68,0.00,0.00}{#1}}
\newcommand{\ExtensionTok}[1]{\textcolor[rgb]{0.00,0.23,0.31}{#1}}
\newcommand{\FloatTok}[1]{\textcolor[rgb]{0.68,0.00,0.00}{#1}}
\newcommand{\FunctionTok}[1]{\textcolor[rgb]{0.28,0.35,0.67}{#1}}
\newcommand{\ImportTok}[1]{\textcolor[rgb]{0.00,0.46,0.62}{#1}}
\newcommand{\InformationTok}[1]{\textcolor[rgb]{0.37,0.37,0.37}{#1}}
\newcommand{\KeywordTok}[1]{\textcolor[rgb]{0.00,0.23,0.31}{\textbf{#1}}}
\newcommand{\NormalTok}[1]{\textcolor[rgb]{0.00,0.23,0.31}{#1}}
\newcommand{\OperatorTok}[1]{\textcolor[rgb]{0.37,0.37,0.37}{#1}}
\newcommand{\OtherTok}[1]{\textcolor[rgb]{0.00,0.23,0.31}{#1}}
\newcommand{\PreprocessorTok}[1]{\textcolor[rgb]{0.68,0.00,0.00}{#1}}
\newcommand{\RegionMarkerTok}[1]{\textcolor[rgb]{0.00,0.23,0.31}{#1}}
\newcommand{\SpecialCharTok}[1]{\textcolor[rgb]{0.37,0.37,0.37}{#1}}
\newcommand{\SpecialStringTok}[1]{\textcolor[rgb]{0.13,0.47,0.30}{#1}}
\newcommand{\StringTok}[1]{\textcolor[rgb]{0.13,0.47,0.30}{#1}}
\newcommand{\VariableTok}[1]{\textcolor[rgb]{0.07,0.07,0.07}{#1}}
\newcommand{\VerbatimStringTok}[1]{\textcolor[rgb]{0.13,0.47,0.30}{#1}}
\newcommand{\WarningTok}[1]{\textcolor[rgb]{0.37,0.37,0.37}{\textit{#1}}}

\usepackage{longtable,booktabs,array}
\newcounter{none} % for unnumbered tables
\usepackage{calc} % for calculating minipage widths
% Correct order of tables after \paragraph or \subparagraph
\usepackage{etoolbox}
\makeatletter
\patchcmd\longtable{\par}{\if@noskipsec\mbox{}\fi\par}{}{}
\makeatother
% Allow footnotes in longtable head/foot
\IfFileExists{footnotehyper.sty}{\usepackage{footnotehyper}}{\usepackage{footnote}}
\makesavenoteenv{longtable}
\usepackage{graphicx}
\makeatletter
\newsavebox\pandoc@box
\newcommand*\pandocbounded[1]{% scales image to fit in text height/width
  \sbox\pandoc@box{#1}%
  \Gscale@div\@tempa{\textheight}{\dimexpr\ht\pandoc@box+\dp\pandoc@box\relax}%
  \Gscale@div\@tempb{\linewidth}{\wd\pandoc@box}%
  \ifdim\@tempb\p@<\@tempa\p@\let\@tempa\@tempb\fi% select the smaller of both
  \ifdim\@tempa\p@<\p@\scalebox{\@tempa}{\usebox\pandoc@box}%
  \else\usebox{\pandoc@box}%
  \fi%
}
% Set default figure placement to htbp
\def\fps@figure{htbp}
\makeatother





\setlength{\emergencystretch}{3em} % prevent overfull lines

\providecommand{\tightlist}{%
  \setlength{\itemsep}{0pt}\setlength{\parskip}{0pt}}



 


\makeatletter
\@ifpackageloaded{tcolorbox}{}{\usepackage[skins,breakable]{tcolorbox}}
\@ifpackageloaded{fontawesome5}{}{\usepackage{fontawesome5}}
\definecolor{quarto-callout-color}{HTML}{909090}
\definecolor{quarto-callout-note-color}{HTML}{0758E5}
\definecolor{quarto-callout-important-color}{HTML}{CC1914}
\definecolor{quarto-callout-warning-color}{HTML}{EB9113}
\definecolor{quarto-callout-tip-color}{HTML}{00A047}
\definecolor{quarto-callout-caution-color}{HTML}{FC5300}
\definecolor{quarto-callout-color-frame}{HTML}{acacac}
\definecolor{quarto-callout-note-color-frame}{HTML}{4582ec}
\definecolor{quarto-callout-important-color-frame}{HTML}{d9534f}
\definecolor{quarto-callout-warning-color-frame}{HTML}{f0ad4e}
\definecolor{quarto-callout-tip-color-frame}{HTML}{02b875}
\definecolor{quarto-callout-caution-color-frame}{HTML}{fd7e14}
\makeatother
\makeatletter
\@ifpackageloaded{caption}{}{\usepackage{caption}}
\AtBeginDocument{%
\ifdefined\contentsname
  \renewcommand*\contentsname{Table of contents}
\else
  \newcommand\contentsname{Table of contents}
\fi
\ifdefined\listfigurename
  \renewcommand*\listfigurename{List of Figures}
\else
  \newcommand\listfigurename{List of Figures}
\fi
\ifdefined\listtablename
  \renewcommand*\listtablename{List of Tables}
\else
  \newcommand\listtablename{List of Tables}
\fi
\ifdefined\figurename
  \renewcommand*\figurename{Figure}
\else
  \newcommand\figurename{Figure}
\fi
\ifdefined\tablename
  \renewcommand*\tablename{Table}
\else
  \newcommand\tablename{Table}
\fi
}
\@ifpackageloaded{float}{}{\usepackage{float}}
\floatstyle{ruled}
\@ifundefined{c@chapter}{\newfloat{codelisting}{h}{lop}}{\newfloat{codelisting}{h}{lop}[chapter]}
\floatname{codelisting}{Listing}
\newcommand*\listoflistings{\listof{codelisting}{List of Listings}}
\makeatother
\makeatletter
\makeatother
\makeatletter
\@ifpackageloaded{caption}{}{\usepackage{caption}}
\@ifpackageloaded{subcaption}{}{\usepackage{subcaption}}
\makeatother
\usepackage{bookmark}
\IfFileExists{xurl.sty}{\usepackage{xurl}}{} % add URL line breaks if available
\urlstyle{same}
\hypersetup{
  pdftitle={Identification of a Forebrain-Enriched Astrocyte Program in the Human Brain Non-Neuronal Atlas},
  pdfauthor={Neuro-AI Brain Cell Atlas Project},
  pdfkeywords={single-cell RNA-seq, astrocyte heterogeneity, Allen Brain
Cell Atlas, neuroprotection},
  colorlinks=true,
  linkcolor={blue},
  filecolor={Maroon},
  citecolor={Blue},
  urlcolor={Blue},
  pdfcreator={LaTeX via pandoc}}


\title{Identification of a Forebrain-Enriched Astrocyte Program in the
Human Brain Non-Neuronal Atlas}
\usepackage{etoolbox}
\makeatletter
\providecommand{\subtitle}[1]{% add subtitle to \maketitle
  \apptocmd{\@title}{\par {\large #1 \par}}{}{}
}
\makeatother
\subtitle{Conference Summary - Neuro-AI Brain Cell Atlas (Notebook 01)}
\author{Neuro-AI Brain Cell Atlas Project}
\date{June 21, 2026}
\begin{document}
\maketitle


\subsection{Conference abstract}\label{conference-abstract}

\textbf{Background.} Astrocytes are heterogeneous across brain regions
and perform distinct roles in synaptic support, ion homeostasis, and
neurovascular coupling. Whether human non-neuronal single-cell atlases
resolve forebrain-specialized astrocyte states remains an open question.

\textbf{Methods.} We analyzed log\textsubscript{2}-normalized
single-nucleus RNA-seq data from the Allen Institute Whole Human Brain
(WHB) atlas (Siletti et al., \emph{Nature} 2023;
\href{https://alleninstitute.github.io/abc_atlas_access/descriptions/WHB-10X-clustering.html}{WHB-10X
clustering}). From 888,263 non-neuronal nuclei, we subsampled 50,000
cells, selected highly variable genes, and computed PCA, a neighbor
graph, UMAP, and Leiden clustering (\texttt{resolution\ =\ 0.5}).
Clusters were manually annotated using canonical markers. Differential
expression, regional composition, and Gene Ontology enrichment (Enrichr)
were performed for the \textbf{Astrocyte\_2} population relative to
canonical astrocytes.

\textbf{Results.} Leiden clustering resolved 12 non-neuronal
populations, including oligodendrocytes, microglia, endothelial cells,
pericytes, and two astrocyte states. \textbf{Astrocyte\_2} was
transcriptionally distinct from \textbf{Astrocyte}, with top markers
including \emph{SLC1A2}, \emph{SPARCL1}, \emph{SLC4A4}, \emph{AQP4},
\emph{APOE}, and \emph{CLU}. Within the Astrocyte\_2 compartment,
\textbf{51.5\%} of cells mapped to cerebral cortex and \textbf{27.9\%}
to hippocampus, with negligible representation in brainstem and spinal
cord. GO enrichment highlighted astrocyte differentiation, long-term
synaptic potentiation, ion transport, neuron projection development, and
regulation of amyloid-beta formation. A composite \textbf{Astrocyte
Support Signature Score} (\emph{APOE}, \emph{CLU}, \emph{GFAP},
\emph{SLC1A2}, \emph{AQP4}, \emph{SPARCL1}, \emph{SLC4A4}) placed
\textbf{96.0\%} of Astrocyte\_2 cells in the highest quartile.

\textbf{Conclusions.} Astrocyte\_2 represents a forebrain-enriched,
synaptic-support astrocyte program consistent with glutamate clearance,
water/ion transport, and neurovascular homeostasis. These findings
motivate follow-up analyses of region-specific astrocyte function and
neurodegenerative disease relevance (Notebook 02).

\begin{center}\rule{0.5\linewidth}{0.5pt}\end{center}

\subsection{Significance}\label{significance}

\begin{tcolorbox}[enhanced jigsaw, arc=.35mm, bottomrule=.15mm, bottomtitle=1mm, breakable, colback=white, colbacktitle=quarto-callout-note-color!10!white, colframe=quarto-callout-note-color-frame, coltitle=black, left=2mm, leftrule=.75mm, opacityback=0, opacitybacktitle=0.6, rightrule=.15mm, title=\textcolor{quarto-callout-note-color}{\faInfo}\hspace{0.5em}{Take-home message}, titlerule=0mm, toprule=.15mm, toptitle=1mm]

Non-neuronal scRNA-seq of the whole human brain resolves at least two
astrocyte states. \textbf{Astrocyte\_2} is concentrated in cortex and
hippocampus and expresses a mature, neuroprotective-support
transcriptional program---not a generic housekeeping glial signature.

\end{tcolorbox}

From a systems-neuroscience perspective, this pattern is biologically
coherent: \emph{SLC1A2} (EAAT2) mediates glutamate uptake at synapses;
\emph{AQP4} and \emph{SLC4A4} support water and bicarbonate homeostasis
at the glial limitans; \emph{SPARCL1} (Hevin) participates in synapse
formation and maturation; and \emph{APOE}/\emph{CLU} link glial biology
to lipid handling and stress response pathways relevant to Alzheimer's
disease genetics.

\begin{center}\rule{0.5\linewidth}{0.5pt}\end{center}

\subsection{Methods summary}\label{methods-summary}

{\def\LTcaptype{none} % do not increment counter
\begin{longtable}[]{@{}
  >{\raggedright\arraybackslash}p{(\linewidth - 2\tabcolsep) * \real{0.4286}}
  >{\raggedright\arraybackslash}p{(\linewidth - 2\tabcolsep) * \real{0.5714}}@{}}
\toprule\noalign{}
\begin{minipage}[b]{\linewidth}\raggedright
Step
\end{minipage} & \begin{minipage}[b]{\linewidth}\raggedright
Detail
\end{minipage} \\
\midrule\noalign{}
\endhead
\bottomrule\noalign{}
\endlastfoot
\textbf{Dataset} & Allen WHB-10Xv3 non-neuron log2 h5ad (888,263 ×
59,357) \\
\textbf{Subset} & Random 50,000 nuclei for compute-efficient
exploration \\
\textbf{Preprocessing} & Highly variable genes → PCA (50 PCs) →
neighbors → UMAP \\
\textbf{Clustering} & Leiden (\texttt{resolution\ =\ 0.5}); 12
clusters \\
\textbf{Annotation} & Manual mapping using marker genes and Allen
metadata \\
\textbf{DE testing} & \texttt{rank\_genes\_groups}: Astrocyte\_2 vs
Astrocyte \\
\textbf{Enrichment} & GSEApy Enrichr - GO Biological Process 2023 \\
\textbf{Scoring} & Scanpy \texttt{score\_genes} - Astrocyte Support gene
set \\
\textbf{Software} & Scanpy, AnnData, GSEApy; conda env
\texttt{brain-cell-atlas} \\
\textbf{Reproducibility} &
\texttt{notebooks/01\_load\_brain\_data.ipynb}; processed object saved
to \texttt{data/processed/brain\_non\_neuronal\_50k\_annotated.h5ad} \\
\end{longtable}
}

\begin{center}\rule{0.5\linewidth}{0.5pt}\end{center}

\subsection{Results}\label{results}

\subsubsection{1. Non-neuronal landscape}\label{non-neuronal-landscape}

Across 14 anatomical divisions represented in the subsample, expected
major glial and vascular populations were recovered after Leiden
clustering and manual curation:

{\def\LTcaptype{none} % do not increment counter
\begin{longtable}[]{@{}ll@{}}
\toprule\noalign{}
Leiden & Annotated cell type \\
\midrule\noalign{}
\endhead
\bottomrule\noalign{}
\endlastfoot
0 & Activated Microglia \\
1 & Oligodendrocyte \\
2, 9 & Unknown \\
3 & Endothelial \\
4 & Astrocyte \\
5 & Pericyte \\
\textbf{6} & \textbf{Astrocyte\_2} \\
7 & Fibroblast \\
8 & Choroid Plexus \\
10 & Ependymal \\
11 & Border Macrophage \\
\end{longtable}
}

\begin{tcolorbox}[enhanced jigsaw, arc=.35mm, bottomrule=.15mm, bottomtitle=1mm, breakable, colback=white, colbacktitle=quarto-callout-tip-color!10!white, colframe=quarto-callout-tip-color-frame, coltitle=black, left=2mm, leftrule=.75mm, opacityback=0, opacitybacktitle=0.6, rightrule=.15mm, title=\textcolor{quarto-callout-tip-color}{\faLightbulb}\hspace{0.5em}{Regional pattern (Research Question 1)}, titlerule=0mm, toprule=.15mm, toptitle=1mm]

Column-normalized region × cell-type tables revealed:

\begin{itemize}
\tightlist
\item
  \textbf{Claustrum} enriched for Activated Microglia
  (\textasciitilde53\%)
\item
  \textbf{Cerebral cortex} enriched for Astrocyte\_2
  (\textasciitilde24\% of cortical non-neurons in subsample)
\item
  \textbf{Basal forebrain} enriched for Choroid Plexus
  (\textasciitilde25\%)
\item
  \textbf{Myelencephalon} enriched for Oligodendrocytes
  (\textasciitilde21\%)
\end{itemize}

\end{tcolorbox}

\subsubsection{2. Astrocyte\_2 is
forebrain-localized}\label{astrocyte_2-is-forebrain-localized}

Among \textbf{Astrocyte\_2} cells, anatomical division composition was:

{\def\LTcaptype{none} % do not increment counter
\begin{longtable}[]{@{}ll@{}}
\toprule\noalign{}
Anatomical division & Fraction of Astrocyte\_2 \\
\midrule\noalign{}
\endhead
\bottomrule\noalign{}
\endlastfoot
Cerebral cortex & \textbf{51.5\%} \\
Hippocampus & \textbf{27.9\%} \\
Amygdaloid complex & 8.9\% \\
Basal nuclei & 3.8\% \\
Hypothalamus & 2.7\% \\
Midbrain / Pons / Brainstem / Spinal cord & \textbf{\textless{} 0.1\%
combined} \\
\end{longtable}
}

This spatial restriction supports a \textbf{forebrain synaptic-support
specialization} rather than a pan-CNS astrocyte program.

\subsubsection{3. Marker genes distinguish Astrocyte\_2 from
Astrocyte}\label{marker-genes-distinguish-astrocyte_2-from-astrocyte}

Top differential genes (Astrocyte\_2 vs Astrocyte reference):

{\def\LTcaptype{none} % do not increment counter
\begin{longtable}[]{@{}llll@{}}
\toprule\noalign{}
Rank & Gene & Score & log\textsubscript{2} FC \\
\midrule\noalign{}
\endhead
\bottomrule\noalign{}
\endlastfoot
1 & \emph{SLC1A2} & 81.7 & 19.7 \\
2 & \emph{ADGRV1} & 81.6 & 19.5 \\
3 & \emph{SPARCL1} & 79.0 & 15.3 \\
4 & \emph{SLC4A4} & 78.7 & 17.1 \\
5 & \emph{WIF1} & --- & --- \\
6 & \emph{EPHA6} & --- & --- \\
\end{longtable}
}

Validation genes with higher expression in Astrocyte\_2: \emph{SLC1A2},
\emph{AQP4}, \emph{SPARCL1}, \emph{SLC4A4}.

\textbf{Biological reading:} Astrocyte\_2 expresses a \textbf{mature,
synapse-facing astrocyte module} - glutamate transport (\emph{SLC1A2}),
extracellular matrix/synapse organization (\emph{SPARCL1}), and
ion/water transport (\emph{AQP4}, \emph{SLC4A4}).

\subsubsection{4. Pathway enrichment supports synaptic and glial
maturation
programs}\label{pathway-enrichment-supports-synaptic-and-glial-maturation-programs}

Top GO Biological Process terms among Astrocyte\_2 marker genes
(Enrichr, adjusted \emph{P} \textless{} 0.05):

{\def\LTcaptype{none} % do not increment counter
\begin{longtable}[]{@{}
  >{\raggedright\arraybackslash}p{(\linewidth - 4\tabcolsep) * \real{0.2143}}
  >{\raggedright\arraybackslash}p{(\linewidth - 4\tabcolsep) * \real{0.3571}}
  >{\raggedright\arraybackslash}p{(\linewidth - 4\tabcolsep) * \real{0.4286}}@{}}
\toprule\noalign{}
\begin{minipage}[b]{\linewidth}\raggedright
Term
\end{minipage} & \begin{minipage}[b]{\linewidth}\raggedright
Adj. \emph{P}
\end{minipage} & \begin{minipage}[b]{\linewidth}\raggedright
Lead genes
\end{minipage} \\
\midrule\noalign{}
\endhead
\bottomrule\noalign{}
\endlastfoot
Regulation of astrocyte differentiation & 0.010 & \emph{EPHA4},
\emph{ID2}, \emph{HES5} \\
Regulation of long-term synaptic potentiation & 0.039 & --- \\
Regulation of amyloid-beta formation & 0.039 & \emph{EPHA4},
\emph{APOE}, \emph{CLU} \\
Negative regulation of long-term synaptic potentiation & 0.042 &
\emph{EPHA4}, \emph{APOE} \\
Positive regulation of neuron projection development & 0.054 &
\emph{NTRK3}, \emph{APOE}, \ldots{} \\
\end{longtable}
}

These pathways align with a glial state engaged in \textbf{circuit
maintenance}, not merely structural support.

\subsubsection{5. Astrocyte Support Signature Score (Research Question
2)}\label{astrocyte-support-signature-score-research-question-2}

\textbf{Question:} Are Astrocyte\_2 cells enriched for Alzheimer's
disease--related / neuroprotective molecular programs?

Because several genes in the scoring panel overlap with the marker genes
used to characterize Astrocyte\_2, this score should be interpreted as a
descriptive summary of the identified transcriptional program rather
than an independent validation.

Gene set scored: \emph{APOE}, \emph{CLU}, \emph{GFAP}, \emph{SLC1A2},
\emph{AQP4}, \emph{SPARCL1}, \emph{SLC4A4}.

{\def\LTcaptype{none} % do not increment counter
\begin{longtable}[]{@{}lllll@{}}
\toprule\noalign{}
Cell type & Q0 & Q1 & Q2 & Q3 (highest) \\
\midrule\noalign{}
\endhead
\bottomrule\noalign{}
\endlastfoot
Astrocyte & 36.9\% & 21.6\% & 28.2\% & 13.3\% \\
\textbf{Astrocyte\_2} & 0.2\% & 0.5\% & 3.3\% & \textbf{96.0\%} \\
Activated Microglia & 44.1\% & 22.6\% & 28.3\% & 5.0\% \\
\end{longtable}
}

\begin{tcolorbox}[enhanced jigsaw, arc=.35mm, bottomrule=.15mm, bottomtitle=1mm, breakable, colback=white, colbacktitle=quarto-callout-important-color!10!white, colframe=quarto-callout-important-color-frame, coltitle=black, left=2mm, leftrule=.75mm, opacityback=0, opacitybacktitle=0.6, rightrule=.15mm, title=\textcolor{quarto-callout-important-color}{\faExclamation}\hspace{0.5em}{Interpretation (professor's conclusion)}, titlerule=0mm, toprule=.15mm, toptitle=1mm]

Astrocyte\_2 is \textbf{not} a generic astrocyte cluster. It is a
\textbf{forebrain-enriched, neuroprotective-support astrocyte state}
with convergent evidence from:

\begin{enumerate}
\def\labelenumi{\arabic{enumi}.}
\tightlist
\item
  Regional anatomy (cortex + hippocampus)
\item
  Synaptic glutamate/ion transport markers
\item
  GO terms for synaptic plasticity and amyloid-beta regulation
\item
  Near-uniform occupancy of the top neuroprotective score quartile
\end{enumerate}

This is \textbf{hypothesis-generating}, not proof of causal
neuroprotection. Subsampling, manual annotation, and absence of spatial
validation are important caveats.

\end{tcolorbox}

\begin{center}\rule{0.5\linewidth}{0.5pt}\end{center}

\subsection{Conclusions}\label{conclusions}

\begin{enumerate}
\def\labelenumi{\arabic{enumi}.}
\item
  \textbf{Single-cell analysis of human non-neuronal nuclei resolves
  astrocyte heterogeneity.} At least two astrocyte states coexist in the
  WHB atlas subsample.
\item
  \textbf{Astrocyte\_2 maps to forebrain cognition circuits.} Strong
  enrichment in cerebral cortex and hippocampus, depletion in
  hindbrain/spinal cord.
\item
  \textbf{Transcriptional identity matches synaptic-support astrocytes.}
  Elevated \emph{SLC1A2}, \emph{AQP4}, \emph{SPARCL1}, \emph{SLC4A4},
  \emph{APOE}, and \emph{CLU}.
\item
  \textbf{Pathway analysis links Astrocyte\_2 to synaptic plasticity and
  amyloid-beta regulation.} Relevant to Alzheimer's disease biology, but
  requires orthogonal validation.
\item
  \textbf{Next experiment (Notebook 02):} Subset hippocampal and
  cortical Astrocyte\_2 cells; test whether neuroprotective gene
  programs covary with region, donor, or disease-associated gene
  modules; compare to Allen taxonomy labels in full metadata.
\end{enumerate}

\begin{center}\rule{0.5\linewidth}{0.5pt}\end{center}

\subsection{Limitations}\label{limitations}

\begin{itemize}
\tightlist
\item
  \textbf{Subsampling:} 50,000 of 888,263 non-neuronal nuclei, rare
  populations may be underrepresented.
\item
  \textbf{Manual cluster labels:} Leiden resolution and annotation
  subject to investigator bias; should be cross-validated against Allen
  WHB taxonomy (\texttt{WHB-taxonomy} metadata).
\item
  \textbf{Log\textsubscript{2} matrix:} Analysis on normalized
  expression; raw-count DE methods not applied.
\item
  \textbf{No spatial context:} Regional labels are dissection-based, not
  in situ MERFISH validation.
\item
  \textbf{Exploratory enrichment:} Enrichr GO results require
  independent cohort replication.
\end{itemize}

\begin{center}\rule{0.5\linewidth}{0.5pt}\end{center}

\subsection{Future directions}\label{future-directions}

\begin{itemize}
\tightlist
\item
  Integrate Allen \textbf{cell\_metadata} + \textbf{WHB-taxonomy} for
  reference-based annotation
\item
  Region-stratified re-clustering (hippocampus vs cortex astrocytes)
\item
  Neuron vs non-neuron comparative atlas (download
  \texttt{WHB-10Xv3-Neurons/log2})
\item
  Spatial validation with Allen MERFISH resources
\item
  Export key figures from Notebook 01 to \texttt{reports/figures/} for
  poster submission
\end{itemize}

\begin{center}\rule{0.5\linewidth}{0.5pt}\end{center}

\subsection{Render this report}\label{render-this-report}

From the project root (Quarto required:
\texttt{conda\ install\ -c\ conda-forge\ quarto}):

\begin{Shaded}
\begin{Highlighting}[]
\ExtensionTok{quarto}\NormalTok{ render reports/brain\_cell\_atlas\_report.qmd}
\end{Highlighting}
\end{Shaded}

Output: \texttt{reports/brain\_cell\_atlas\_report.html} (and PDF if
LaTeX installed).

\begin{center}\rule{0.5\linewidth}{0.5pt}\end{center}

\subsection{References}\label{references}

\textbf{Primary data}

\begin{enumerate}
\def\labelenumi{\arabic{enumi}.}
\tightlist
\item
  Siletti K, et al.~Transcriptomic diversity of cell types across the
  adult human brain. \emph{Nature}. 2023.
\item
  Allen Institute for Brain Science.
  \href{https://alleninstitute.github.io/abc_atlas_access/descriptions/WHB-10X-clustering.html}{WHB-10X
  clustering} and
  \href{https://alleninstitute.github.io/abc_atlas_access/descriptions/WHB-10Xv3.html}{WHB-10Xv3
  expression}. \emph{abc\_atlas\_access} documentation.
\end{enumerate}

\textbf{Software}

\begin{enumerate}
\def\labelenumi{\arabic{enumi}.}
\setcounter{enumi}{2}
\tightlist
\item
  Wolf FA, Angerer P, Theis FJ. SCANPY: large-scale single-cell gene
  expression data analysis. \emph{Genome Biology}. 2018.
\item
  Fang Z, et al.~GSEApy: a comprehensive package for performing gene set
  enrichment analysis in Python. \emph{Bioinformatics}. 2023.
\end{enumerate}

\textbf{Key genes (quick reference)}

{\def\LTcaptype{none} % do not increment counter
\begin{longtable}[]{@{}
  >{\raggedright\arraybackslash}p{(\linewidth - 2\tabcolsep) * \real{0.5000}}
  >{\raggedright\arraybackslash}p{(\linewidth - 2\tabcolsep) * \real{0.5000}}@{}}
\toprule\noalign{}
\begin{minipage}[b]{\linewidth}\raggedright
Gene
\end{minipage} & \begin{minipage}[b]{\linewidth}\raggedright
Role
\end{minipage} \\
\midrule\noalign{}
\endhead
\bottomrule\noalign{}
\endlastfoot
\emph{SLC1A2} & Glutamate transporter EAAT2 - synaptic clearance \\
\emph{AQP4} & Aquaporin-4 - astrocytic water homeostasis, glymphatic
interface \\
\emph{SPARCL1} & Synapse-organizing glycoprotein (Hevin) \\
\emph{SLC4A4} & Bicarbonate transporter - pH / ion homeostasis \\
\emph{APOE} & Lipid transport; major AD risk locus \\
\emph{CLU} & Clusterin; chaperone / stress response; AD-associated \\
\end{longtable}
}




\end{document}
